% -*- coding: utf-8 -*-
% !TEX program = xelatex
\documentclass[plain,noanswer,medmath]{../../template/jnuexam}
\usepackage{../../template/kysx}

\begin{document}


\examtitle{2007}{1}


\makepart{选择题}{1～10小题，每小题4分，共40分}


\begin{problem}
当$x \to 0^+$时，与$\sqrt{x}$等价的无穷小量是\pickout{}
\begin{abcd}
\item $1 - \e^{\sqrt{x}}$．
\item $\ln \frac{1 + x}{1 - \sqrt{x}}$．
\item $\sqrt{1 + \sqrt{x}} - 1$．
\item $1 - \cos \sqrt{x}$．
\end{abcd}
\end{problem}


\begin{problem}
曲线$y = \frac{1}{x} + \ln(1 + \e^x)$渐近线的条数为\pickout{}
\begin{abcd}
\item $0$．
\item $1$．
\item $2$．
\item $3$．
\end{abcd}
\end{problem}


\begin{problem}
如图，连续函数$y = f(x)$在区间$[- 3, - 2]$, $[2,3]$上的图形分别是直径为$1$的上、下半圆周，在区间$[- 2,0]$, $[0,2]$上图形分别是直径为$2$的上、下半圆周．
\par\centerline{\includegraphics[width=0.5\linewidth]{2007A1}}
设$F(x) = \int_0^x f(t) \dt$，则下列结论正确的是\pickout{}
\begin{abcd}
\item $F(3) = - \frac{3}{4}F(- 2)$．
\item $F(3) = \frac{5}{4}F(2)$．
\item $F(- 3) = \frac{3}{4}F(2)$．
\item $F(- 3) = - \frac{5}{4}F(- 2)$．
\end{abcd}
\end{problem}


\begin{problem}
设函数$f(x)$在$x = 0$连续，则下列命题错误的是\pickout{}
\begin{abcd}
\item 若$\lim_{x \to 0} \frac{f(x)}{x}$存在，则$f(0) = 0$．
\item 若$\lim_{x \to 0} \frac{f(x) + f(- x)}{x}$存在，则$f(0) = 0$．
\item 若$\lim_{x \to 0} \frac{f(x)}{x}$存在，则$f'(0)$存在．
\item 若$\lim_{x \to 0} \frac{f(x) - f(- x)}{x}$存在，则$f'(0)$存在．
\end{abcd}
\end{problem}


\begin{problem}
设函数$f(x)$在$(0, +\infty)$上具有二阶导数，且$f''(x) > 0$，令$u_n = f(n)$（$n=1,2,\cdots$），
则下列结论正确的是\pickout{}
\begin{abcd}
\item 若$u_1 > u_2$，则$\left\{u_n \right\}$必收敛．
\item 若$u_1 > u_2$，则$\left\{u_n \right\}$必发散．
\item 若$u_1 < u_2$，则$\left\{u_n \right\}$必收敛．
\item 若$u_1 < u_2$，则$\left\{u_n \right\}$必发散．
\end{abcd}
\end{problem}


\begin{problem}
设曲线$L:f(x,y) = 1$（$f(x,y)$具有一阶连续偏导数）过第Ⅱ象限内的点$M$和第Ⅳ象限内的点$N$，
$\Gamma$为$L$上从点$M$到点$N$的一段弧，则下列积分小于零的是\pickout{}
\begin{abcd}
\item $\int_{\Gamma}f(x,y)\dx$．
\item $\int_{\Gamma}f(x,y)\dy$．
\item $\int_{\Gamma}f(x,y)\ds$．
\item $\int_{\Gamma}f'_x(x,y)\dx + f'_y(x,y)\dy$．
\end{abcd}
\end{problem}


\begin{problem}
设向量组$\alpha_1$, $\alpha_2$, $\alpha_3$线性无关，则下列向量组线性相关的是\pickout{}
\begin{abcd}
\item $\alpha_1 - \alpha_2$, $\alpha_2 - \alpha_3$, $\alpha_3 - \alpha_1$．
\item $\alpha_1 + \alpha_2$, $\alpha_2 + \alpha_3$, $\alpha_3 + \alpha_1$．
\item $\alpha_1 - 2\alpha_2$, $\alpha_2 - 2\alpha_3$, $\alpha_3 - 2\alpha_1$．
\item $\alpha_1 + 2\alpha_2$, $\alpha_2 + 2\alpha_3$, $\alpha_3 + 2\alpha_1$．
\end{abcd}
\end{problem}


\begin{problem}
设矩阵$A = \begin{pmatrix}
   2 & -1 & -1 \\
  -1 &  2 & -1 \\
  -1 & -1 & 2
\end{pmatrix}$，$B = \begin{pmatrix}
  1 & 0 & 0 \\
  0 & 1 & 0 \\
  0 & 0 & 0
\end{pmatrix}$，则$A$与$B$\pickout{}
\begin{abcd}
\item 合同，且相似．
\item 合同，但不相似．
\item 不合同，但相似．
\item 既不合同，也不相似．
\end{abcd}
\end{problem}


\begin{problem}
某人向同一目标独立重复射击，每次射击命中目标的概率为$p$（$0 < p < 1$），
则此人第$4$次射击恰好第$2$次命中目标的概率为 \pickout{}
\begin{abcd}
\item $3p(1 - p)^2$．
\item $6p(1 - p)^2$．
\item $3p^2(1 - p)^2$．
\item $6p^2(1 - p)^2$．
\end{abcd}
\end{problem}


\begin{problem}
设随机变量$(X,Y)$服从二维正态分布，且$X$与$Y$不相关，$f_X(x),f_Y(y)$分别表示$X,Y$的概率密度，
则在$Y = y$条件下，$X$的条件概率密度$f_{X\left| Y \right.}(x\left| y) \right.$为\pickout{}
\begin{abcd}
\item $f_X(x)$．
\item $f_Y(y$．
\item $f_X(x) f_Y(y)$．
\item $\frac{f_X(x)}{f_Y(y)}$．
\end{abcd}
\end{problem}


\makepart{填空题}{11～16小题，每小题4分，共24分}


\begin{problem}
$\int_1^2\frac{1}{x^3} \e^{\frac{1}{x}} \dx =$\fillin{}．
\end{problem}


\begin{problem}
设$f(u,v)$为二元可微函数，$z = f(x^y,y^x)$，则$\frac{\pdz}{\pdx} =$ \fillin{}．
\end{problem}


\begin{problem}
二阶常系数非齐次线性微分方程$y'' - 4y' + 3y = 2\e^{2x}$的通解为$y=$ \fillin{}．
\end{problem}


\begin{problem}
设曲面$\Sigma :|x| + |y| + |z| = 1$，
则$\oiint_{\Sigma}(x + |y|) \dS =$ \fillin{}．
\end{problem}


\begin{problem}
设矩阵$A = \begin{pmatrix}
  0 & 1 & 0 & 0 \\
  0 & 0 & 1 & 0 \\
  0 & 0 & 0 & 1 \\
  0 & 0 & 0 & 0
\end{pmatrix}$，则$A^3$的秩为\fillin{}．
\end{problem}


\begin{problem}
在区间$(0,1)$中随机地取两个数，则这两数之差的绝对值小于$\frac{1}{2}$的概率为\fillin{}．
\end{problem}


\makepart{解答题}{17～24小题，共86分}


\begin{problem}[points=10]
求函数$f(x,y) = x^2 + 2y^2 - x^2y^2$在区域
$D = \left\{(x,y)\left| x^2 + y^2 \le 4,y \ge 0 \right. \right\}$上的最大值和最小值．
\end{problem}


\begin{problem}[points=11]
计算曲面积分$I = \iint_{\Sigma}xz\dy\dz + 2zy\dz\dx + 3xy \dx\dy$，
其中$\Sigma$为曲面$z = 1 - x^2 - \frac{y^2}{4}$（$0 \le z \le 1$）的上侧．
\end{problem}


\begin{problem}[points=11]
设函数$f(x)$，$g(x)$在$\;\left[a,b \right]$上连续，在$(a,b)$内二阶可导且存在相等的最大值，
又$f(a)$=$g(a)$，$f(b)$=$g(b)$，证明：存在$\xi \in(a,b)$，使得$f''(\xi) = g''(\xi)$．
\end{problem}


\begin{problem}[points=10]
设幂级数$\sum_{n = 0}^{\infty}a_n x^n$在$(-\infty , +\infty)$内收敛，其和函数$y(x)$满足
$$y'' - 2xy' - 4y = 0,y(0) = 0,y'(0) = 1.$$
\begin{enumerate*}
  \item 证明$a_{n + 2} = \frac{2}{n + 1}a_n,\;n = 1,2, \cdots$；
  \item 求$y(x)$的表达式．
\end{enumerate*}
\end{problem}

\begin{problem}[points=11]
设线性方程组①$\begin{cases}
  x_1 + x_2 + x_3 = 0 \\
  x_1 + 2x_2 + ax_3 = 0 \\
  x_1 + 4x_2 + a^2x_3 = 0
\end{cases}$
与方程②       $x_1 + 2x_2 + x_3 = a - 1$
有公共解，求$a$的值及所有公共解．
\end{problem}


\begin{problem}[points=11]
设$3$阶实对称矩阵$A$的特征值$\lambda_1 = 1$, $\lambda_2 = 2$, $\lambda_3 = - 2$，
$\alpha_1 = (1, - 1,1)^T$是$A$的属于$\lambda_1$的一个特征向量，
记$B = A^5 - 4A^3 + E$，其中$E$为$3$阶单位矩阵．
\begin{enumerate}
  \item 验证$\alpha_1$是矩阵$B$的特征向量，并求$B$的全部特征值与特征向量；
  \item 求矩阵$B$．
\end{enumerate}
\end{problem}


\begin{problem}[points=11]
设二维随机变量$(X,Y)$的概率密度为
$$f(x,y) = \begin{cases}
  2 - x - y, & 0 < x < 1,0 < y < 1, \\
  0, & \text{其他}.
\end{cases}$$
\begin{enumerate*}
  \item 求$P\{X > 2Y\}$；
  \item 求$Z = X + Y$的概率密度$f_Z(z)$．
\end{enumerate*}
\end{problem}

\begin{problem}[points=11]
设总体$X$的概率密度为
$$f(x;\theta) = \begin{cases}
  \frac{1}{2\theta}, & 0 < x < \theta , \\
  \frac{1}{2(1-\theta)}, & \theta \le x < 1, \\
  0, & \text{其他}.
\end{cases}$$
其中参数$\theta$（$0 < \theta < 1$）未知，$X_1,X_2,\cdots X_n$是来自总体$X$的简单随机样本，
$\overline{X}$是样本均值．
\begin{enumerate}
  \item 求参数$\theta$的矩估计量$\widehat{\theta}$；
  \item 判断$4\overline{X}^2$是否为$\theta^2$的无偏估计量，并说明理由．
\end{enumerate}
\end{problem}


\end{document}
